\chapter{TURP (Transurethral resection of the prostate)}
\index{TURP (Transurethral resection of the prostate)}

\medskip
\noindent\textit{\textbf{Under review.} This chapter is an AI-assisted educational draft; it carries no source citations and is not a substitute for professional medical advice.}

\section*{Definition and Overview}
TURP (Transurethral resection of the prostate) is a surgical procedure performed to diagnose, treat, correct, or palliate a medical condition. Surgery is one of the principal therapeutic modalities in medicine and may be classified as elective, urgent, or emergency; open or minimally invasive (keyhole); and curative, palliative, or diagnostic. All surgical interventions carry inherent risks, and the decision to operate requires careful consideration of benefits, alternatives, and patient fitness.

\section*{Epidemiology}
Hundreds of millions of surgical procedures are performed globally each year. Surgical volume and complexity have increased with ageing populations and advances in anaesthesia and perioperative care. Postoperative complications occur in approximately 10--15\% of procedures, with major complications in 2--4\%. Surgical safety is improved by standardised checklists, team communication, and protocol-driven care.

\section*{Aetiology and Pathophysiology}
Surgical intervention is indicated for:

\begin{itemize}
    \item Removal of diseased or neoplastic tissue
    \item Repair or reconstruction of structures damaged by injury, congenital anomaly, or degeneration
    \item Relief of mechanical obstruction or compression
    \item Correction of anatomical abnormalities
    \item Diagnostic biopsy or staging
    \item Implantation of prosthetic devices or donor organs
    \item Palliation of incurable conditions (e.g.\ bypass of malignant obstruction)
\end{itemize}

\section*{Clinical Features}
\begin{itemize}
    \item Pre-operative symptoms relate to the underlying condition necessitating surgery
    \item Post-operative: pain and swelling at the operative site
    \item Fatigue and reduced mobility for days to weeks
    \item Mild fever in the first 24--48 hours is common
    \item Reduced appetite and altered bowel function (especially after abdominal surgery)
    \item Wound drainage -- serous or blood-tinged fluid is normal in the early period
\end{itemize}

\section*{Differential Diagnosis}
Postoperative complications to distinguish from normal recovery include:

\begin{itemize}
    \item Wound infection vs.\ normal inflammatory response
    \item Deep vein thrombosis vs.\ postoperative calf pain
    \item Pulmonary embolism vs.\ postoperative dyspnoea
    \item Atelectasis vs.\ pneumonia
    \item Ileus vs.\ mechanical bowel obstruction
    \item Haematoma vs.\ seroma vs.\ wound infection
    \item Urinary retention vs.\ urinary tract infection
\end{itemize}

\section*{When to Seek Medical Care}
After surgery, contact the surgical team or seek urgent care if:
\begin{itemize}
    \item Severe or worsening pain not controlled by prescribed analgesia
    \item Heavy or persistent bleeding from the wound
    \item Spreading redness, warmth, swelling, or purulent discharge (wound infection)
    \item Fever above 38 degrees Celsius that does not settle
    \item Chest pain, dyspnoea, or calf pain and swelling (possible thromboembolism)
    \item Inability to pass urine or open the bowels
    \item Wound dehiscence (separation of wound edges)
\end{itemize}

\textbf{Contact your local emergency services} for any life-threatening postoperative emergency.

\section*{Diagnosis and Investigations}
\begin{itemize}
    \item \textbf{Pre-operative assessment:} full history, examination, and medication review
    \item \textbf{Blood tests:} full blood count, coagulation, electrolytes, group and save
    \item \textbf{Cardiac assessment:} ECG; echocardiogram or cardiology review for cardiac patients
    \item \textbf{Imaging:} X-ray, ultrasound, CT, or MRI to plan the procedure
    \item \textbf{Anaesthetic assessment:} airway, fitness grading (ASA classification), and discussion of anaesthetic options
    \item \textbf{Informed consent:} discussion of risks, benefits, alternatives, and expected recovery
\end{itemize}

\section*{Management}
\begin{itemize}
    \item \textbf{Anaesthesia:} local, regional (e.g.\ spinal, epidural), or general depending on the procedure
    \item \textbf{Operative technique:} open or minimally invasive (laparoscopic, robotic, endoscopic)
    \item \textbf{Perioperative care:} prophylactic antibiotics, VTE prophylaxis, temperature and fluid management
    \item \textbf{Pain management:} multimodal analgesia (paracetamol, NSAIDs, opioids, regional techniques)
    \item \textbf{Wound care:} sutures, clips, dressings; removal at appropriate interval
    \item \textbf{Rehabilitation:} physiotherapy, graded return to activity, and functional restoration
    \item \textbf{Follow-up:} review to assess healing, histopathology results, and ongoing management
\end{itemize}

\section*{Complications}
\begin{itemize}
    \item \textbf{Infection:} wound, deep, or organ/space surgical site infection
    \item \textbf{Bleeding:} intra-operative or postoperative haemorrhage requiring transfusion or re-operation
    \item \textbf{Thromboembolism:} deep vein thrombosis and pulmonary embolism
    \item \textbf{Anaesthetic complications:} aspiration, anaphylaxis, malignant hyperthermia (rare)
    \item \textbf{Organ injury:} inadvertent damage to adjacent structures
    \item \textbf{Wound dehiscence and hernia:} failure of wound healing
    \item \textbf{Systemic:} ileus, urinary retention, pneumonia, acute kidney injury
\end{itemize}

\section*{Prognosis and Outcomes}
Most elective surgical procedures have excellent outcomes, with the majority of patients returning to normal function within weeks to months. Recovery time depends on the procedure, patient factors (age, co-morbidity, fitness), and the presence of complications. Minimally invasive techniques generally result in shorter hospital stays and faster recovery than open surgery. Major complications, when they occur, can significantly prolong recovery and increase mortality.

\section*{Prevention and Self-Care}
\begin{itemize}
    \item Stop smoking well before surgery to reduce wound and anaesthetic complications
    \item Maintain a healthy weight and optimise chronic conditions (diabetes, hypertension) pre-operatively
    \item Follow fasting and medication instructions precisely
    \item Mobilise early postoperatively to reduce DVT and pneumonia risk
    \item Perform prescribed breathing exercises and physiotherapy
    \item Attend all follow-up appointments
    \item Report any concerning symptoms promptly
\end{itemize}
